\begin{abstract}
Operating System policies are becoming increasingly workload-dependent, where assumptions about workloads, application behavior, and hardware shape policy design. When these assumptions fail to hold in deployment, policies can silently degrade performance.
Existing engineering practice relies on offline benchmarking, which is ill-suited to addressing deployment-time performance regressions and provides little insight for future policy iterations. 
This paper argues that OSes need \emph{performance introspection}: the ability to monitor \emph{during execution} whether a policy is satisfying its performance requirements, and---crucially---to intervene when it is not.
To this end we propose the abstraction of \emph{OS guardrails}, which pair (i) a runtime-checkable performance condition with (ii) a corrective action.
We present principles for designing useful guardrails, a Guardrail Description Language (GDL) and semantics for specifying guardails over kernel-visible signals, an implementation substrate based on in-kernel monitors and eBPF, and a verification interface for verifying that local guardrail conditions (and their corrective actions) ensure system-wide performance objectives.
We demonstrate the effectiveness of this approach through five case studies, showing that guardrails expose broken deployment-time assumptions, surface policy-specific failure modes, and enable targeted corrections or safer fallbacks.
\end{abstract}